% PDF 5
\sourcepage{1}{5}
\section*{一个世纪以前的代数几何！}\addcontentsline{toc}{section}{一个世纪以前的代数几何！}
下面是克莱因的《论黎曼的代数函数及其积分理论》（\textit{On Riemann's Theory of Algebraic Functions and their Integrals}，1880；Dover，1963）目录及正文首页的翻译。
\begin{center}
\begin{tabular}{p{.82\linewidth}r}
\toprule
\textbf{第一部分}&页码\\
1. 以平面定常流解释 $x+iy$ 的函数&1\\
2. 考察 $w=f(z)$ 的无穷大&4\\
3. 有理函数及其积分：由低阶无穷大导出高阶无穷大&9\\
4. 用实验产生这些流动&12\\
5. 转到球面：任意曲面上的流动&15\\
6. 前述理论与复变量函数之间的联系&19\\
7. 再论球面流动：黎曼的一般问题&21\\
\midrule
\textbf{第二部分：黎曼的理论}&\\
8. 按整数 $p$ 的值对闭曲面分类&23\\
9. 初步确定任意曲面上的定常流&26\\
10. 最一般的定常流：证明其他流动不可能存在&29\\
11. 通过实例说明这些流动&32\\
12. 由各个分项合成最一般的位置函数&37\\
\bottomrule
\end{tabular}
\end{center}
\subsection*{§1. 以平面定常流解释 $x+iy$ 的函数}
以下各页讨论的 $x+iy$ 的那些函数，其物理解释是众所周知的。\footnote{尤其可参阅麦克斯韦的《电磁学通论》（剑桥，1873）。就对问题的直观处理而言，他的观点与本文所采用的完全相同。}这里仅为完整起见，简述这种解释所依据的原理。
设 $w=u+iv$、$z=x+iy$、$w=f(z)$。首先有
\[
\frac{\partial u}{\partial x}=\frac{\partial v}{\partial y},\qquad
\frac{\partial u}{\partial y}=-\frac{\partial v}{\partial x}.\tag{1}
\]
从而
\[
\frac{\partial^2u}{\partial x^2}+\frac{\partial^2u}{\partial y^2}=0,\tag{2}
\]
对 $v$ 也有
\[
\frac{\partial^2v}{\partial x^2}+\frac{\partial^2v}{\partial y^2}=0.\tag{3}
\]
有关准确的历史细节，请参阅迪厄多内的 D-HISTORY 一文（见参考文献）。

% PDF 6
\sourcepage{2}{6}
\section*{参考文献}\addcontentsline{toc}{section}{参考文献}
下列条目同时列出本书使用的引文简称。为便于查找原作，保留原文书名，并附中文译名。
\begin{description}[style=nextline,leftmargin=1.5em]
\item[ATIYAH--MACDONALD]
Atiyah--Macdonald，\textit{Introduction to Commutative Algebra}（《交换代数导论》），Addison-Wesley，1969。
\item[BOURBAKI]
N. Bourbaki，\textit{Éléments de Mathématique}（《数学原本》），\textit{Actualités scientifiques et industrielles}，Hermann，巴黎。其中：
\begin{itemize}
\item \textit{Algèbre linéaire}（线性代数），第3章，第1236号；简称 BOURBAKI-LINEAR ALGEBRA。
\item \textit{Algèbre multilinéaire}（多重线性代数），第3章，第1044号；简称 BOURBAKI-MULTILINEAR ALGEBRA。
\item \textit{Algèbre commutative}（交换代数），第1290、1293、1308、1314号；简称 BOURBAKI-COMMUTATIVE ALGEBRA。
\end{itemize}
\item[CARTAN--EILENBERG]
H. Cartan、S. Eilenberg，\textit{Homological Algebra}（《同调代数》），Princeton Mathematical Series，Princeton University Press，1956。
\item[DEMAZURE--GABRIEL]
M. Demazure、P. Gabriel，\textit{Groupes algébriques}（《代数群》），North-Holland，阿姆斯特丹。
\item[DIEUDONNÉ]
\textit{Advances in Mathematics}，第3卷第3期，1969年7月：\textit{Algebraic Geometry}（《代数几何》），第233页起，简称 D-ADVANCES, I；\textit{Fondements de la Géométrie Algébrique}（《代数几何基础》），第322页起，简称 D-ADVANCES, II。
另有 \textit{The Historical Development of Algebraic Geometry}（《代数几何的历史发展》），\textit{American Mathematical Monthly}，第79卷第8期，第827--866页；简称 D-HISTORY。
\item[GODEMENT]
R. Godement，\textit{Théorie des faisceaux}（《层论》），\textit{Actualités scientifiques et industrielles} 第1252号，Hermann，巴黎，1958。
\item[GROTHENDIECK：布法罗讲义]
\textit{Survey on the functorial approach to affine algebraic groups}（《仿射代数群的函子方法概述》），布法罗暑期课程，1973；简称 AAG-BUFFALO NOTES。
\textit{Lectures on Topoi}（《拓扑斯讲义》），布法罗暑期课程，1973；简称 TOPOI-BUFFALO NOTES。
\item[TOHOKU]
A. Grothendieck，\textit{Sur quelques points d'algèbre homologique}（《同调代数的若干问题》），\textit{Tohoku Mathematical Journal}，第IX卷（1957），第119--221页。
\item[FONDEMENTS]
A. Grothendieck，\textit{Fondements de la géométrie algébrique}（《代数几何基础》），摘自1957--1962年布尔巴基讨论班；巴黎，Secrétariat mathématique，11, rue Pierre Curie, 5e。
\item[EGA；EGA-I-Springer]
Grothendieck--Dieudonné，\textit{Éléments de Géométrie Algébrique}（《代数几何原理》），\textit{Publications mathématiques de l'Institut des Hautes Études Scientifiques}（IHÉS）；第一卷 Springer 版为第166卷，简称 EGA-I-Springer。
\item[SGA-1-2-3-\ldots]
\textit{Séminaire de géométrie algébrique}（《代数几何讨论班》）。除下述 SGA-2 外，SGA 各卷均重印于 Springer 的 \textit{Lecture Notes in Mathematics}（LNM）丛书：
\begin{itemize}
\item SGA-2：\textit{Cohomologie locale des faisceaux cohérents et Théorèmes de Lefschetz locaux et globaux}（《凝聚层的局部上同调与局部及整体莱夫谢茨定理》），\textit{Advanced Studies in Pure Mathematics}，Masson et Cie 与 North-Holland。
\item SGA-1：\textit{Revêtements étales et Groupe fondamental}（《平展覆叠与基本群》），LNM 224。
\item SGA-3（I、II、III）：\textit{Schémas en groupes}（《群概形》），LNM 151、152、153。
\item SGA-4（I、II、III）：\textit{Théorie des Topos et Cohomologie Étale des schémas}（《拓扑斯理论与概形的平展上同调》），LNM 269、270、305。
\end{itemize}
\end{description}

% PDF 7
\sourcepage{3}{7}
\begin{description}[style=nextline,leftmargin=1.5em]
\item[KÄHLER]
E. Kähler，\textit{Geometria arithmetica}（《算术几何》），\textit{Annali di Matematica}，第4辑，第XLV卷（1958），第1--368页。
\item[MACDONALD-Schemes]
Macdonald，\textit{Algebraic Geometry, Introduction to Schemes}（《代数几何：概形导论》），Benjamin，1968。
\item[MACLANE-Homology；MACLANE-Category text]
S. Mac Lane，\textit{Homology}（《同调》），Springer，柏林--哥廷根--海德堡，1963；\textit{Categories for the Working Mathematician}（《工作数学家的范畴论》），Springer 的 \textit{Graduate Texts in Mathematics} 丛书。
\item[LANG]
Lang，\textit{Algebra}（《代数学》），Addison-Wesley，1968。
\item[MANIN-（{\cyrillicfont МГУ}）]
Y. Manin，\textit{Lectures on Algebraic Geometry}（《代数几何讲义》），莫斯科国立大学（{\cyrillicfont МГУ}）。
\item[MUMFORD]
Mumford，\textit{Introduction to Algebraic Geometry}（《代数几何导论》），哈佛大学油印讲义，未注日期。
\item[FAC]
J.-P. Serre，\textit{Faisceaux algébriques cohérents}（《凝聚代数层》），\textit{Annals of Mathematics}，第LXI卷（1955），第197--278页。
\item[COH-SCHEMES]
\textit{Dix exposés sur la cohomologie des schémas}（《概形上同调十讲》），J. Giraud、A. Grothendieck、S. L. Kleiman、M. Raynaud、J. Tate 著；North-Holland，阿姆斯特丹；Masson，巴黎，1968。
\end{description}

\BookChapter{0}{绪论}
急于进入正题的读者可以先读本章第1节，再读第一章第3节，依此往下；前面的各节，留待需要时再读也无妨。

\section*{前言}\addcontentsline{toc}{section}{前言}
这份讲义主要根据格罗滕迪克1973年夏季访问纽约州立大学布法罗分校时的授课录音整理而成。我又在录音内容之外，补充了习题、与经典代数几何有关的参考材料、历史评述，以及 SGA、EGA 等“圣经”式著作中的具体引文。\footnote{正文中的方括号标记表示由我（F. Gaeta）补入的内容。}\footnote{引号中的作者姓名、书名或论文名等，均指参考文献中的相应条目。}

格罗滕迪克本人不对这份讲义的出版承担任何责任。不过，我相信：目前既然还没有合适的“教科书”，而原始文献对初学者又有相当大的难度，这样的出版物将能帮助更多的读者。本讲义旨在引导读者阅读 SGA、EGA 等原始文献，并具体标明章、节和页码。我还在参考文献中列入迪厄多内的文章、芒福德的讲义等其他入门材料。它们多数只给出证明梗概，甚至完全不证明；或者面向另一类读者，例如 MACDONALD-Schemes 是为经典代数几何学家写的。我希望这份主要面向初入门研究生的讲义，能在这些入门材料与原始文献之间架起桥梁。为帮助新人，这里给出的细节比此类非正式出版物的惯例多得多。我利用口头讲授中某些重复的内容，

% PDF 8
\sourcepage{4}{8}
在大多数节的开头加入了“提要”：主要取自授课录音；如果找不到更好的材料，就由我自己拟写。许多证明是完整的，另一些也几乎完整，只把极少数确实很简单的细节留给读者。

虽然格罗滕迪克本人举过平面代数曲线、曲面等例子，这门课并不预设读者懂得“旧式”或“经典”代数几何。我在不少地方，尤其是为有志于应用数学者写的引言和课程提要中，尝试与以代数簇而非概形为研究对象的“旧式”代数几何建立联系。这么做或许看似背离格罗滕迪克的数学精神，却绝不背离他当时在哲学与社会问题上的关切。在此前那次布法罗之行以及许多其他场合，他都曾反对专家知识和技术。我们怎能忽略这样的事实：不少人翻开一本代数几何书，如果第一页没有出现“代数曲线”或“代数曲面”，就会失望；还有人仅凭“听说”，事先便抱怨代数和范畴语言太多，几何究竟在哪里？我试着消除这些心理上的障碍和成见，借此突出格罗滕迪克所带来的重大简化。为应用数学者写的引言，面向任何具有数学学士学位的人；理论物理学家和工程师应当也能读懂。

我希望，在整门课程最后修订完毕后，讨论概形范畴的第二部分能很快问世。

许多到场听课的同事和学生协助我整理了讲义，订正了排印、语言、数学等各方面的错误。我感谢他们，尤其是 J. Duskin、B. Fell、L. Gupta、P. Hamsher、N. Kazarinoff、M. Klun、I. Ozaki、F. C. Lin、S. Schanuel、G. Sicherman 和 J. Winthrop。首先，我特别感谢格罗滕迪克：他待人如此亲切，又如此慷慨地投入时间。他好几次连续讲了将近七小时，中间只短暂休息几次。谁还能相信他已经对数学失去兴趣了呢？

最后，我同样衷心感谢打字员盖尔·贝尔蒂女士。她出色地完成了工作，以天使般的耐心，几十次订正和重打手稿，却一次也没有抱怨。
\begin{flushright}
费德里科·加埃塔\\
布法罗，1974年6月
\end{flushright}

\Section{0}{向应用数学者推荐代数几何}
\footnote{撰写本节时，M. Barry Fell 提供了许多宝贵建议，我尤其感谢他。}
不过一个世纪以前，纯数学与应用数学之间的区分，在很大程度上还是人为的，也并不重要。例如，克莱因的小书《论黎曼的代数函数及其积分理论》（1880；Dover，1963；参见前面重录的目录），就是从平面或任意曲面上的流体层流这一实际物理问题，来引入黎曼曲面的研究。他甚至在第一页就引用了麦克斯韦的著作！

在我们这个时代，这条“超越方法”的自然延续，是通过霍奇型紧凯勒流形发展起来的复代数流形研究。尽管这种分析方法为进入代数几何这座大厦提供了最美的“入口”之一，人们并不认为它是最根本的入口。在马克斯·诺特之后（约1870年），多年来的主要“入口”是真正的“代数几何”方法：最初它与代数簇的特殊射影模型研究相联系，因而也与不变量论相联系。布里尔和诺特发表于《数学年刊》的经典论文（1874），从双有理的观点奠定了“代数曲线上的几何”的基础。\footnote{译注：原书本段将1870、1874分别误印为1970、1974；按上下文订正。}\footnote{两个不可约代数簇双有理等价，当且仅当它们的有理函数域同构。经典代数几何主要研究不可约代数簇的双有理等价类。}在这种方法中，应用性和具体性依然清清楚楚，从未受到质疑。后来，代数，尤其是交换代数的影响逐渐增强，研究也越来越抽象；论题的表述则越来越脱离应用。

% PDF 9
\sourcepage{5}{9}
今天，在许多同行眼中，格罗滕迪克的代数几何仿佛是当代数学思想最抽象、最无从应用的产物之一。早在美国数学学生开始为学术职位短缺而担忧以前，这种成见就已经造成了损害。如果他们又听到格罗滕迪克在“生存”演讲中批判现代科学、研究、技术等等，恐怕就会更加不安。我记得，当他问听众这些东西有什么用时，我举过那个经典的例子：如果希腊人从未研究圆锥曲线，开普勒又怎能提出他的天体力学定律？当时也谈到了电磁学，并讨论了它可能造成的有害后果。他是个非常开明的人；尽管有这些看法，他仍允许我们使用许多录音机！

我们想说明的是：格罗滕迪克原来的讲述尽管看起来非常抽象，也很少涉及可能的应用，其灵感却十分具体。

与代数几何相比，人们对微分几何的普遍看法完全不同。微分几何从未失去它的应用色彩。作为一种推荐，我希望在这篇引言中说明：像微分流形这样实用的结构，本来就是局部赋环空间的自然例子。因此，熟悉微分流形的读者，就不该觉得格罗滕迪克的概形抽象得那么可怕。诚然，本课程并不把微分流形知识作为逻辑上的先修要求；但关心应用的学生，也应当关心微分流形。

这篇非正式引言的目的，是在格罗滕迪克引入的新数学对象与数学物理中的某些对象之间建立类比。

我选一个自己感兴趣的应用。考虑一个具有 $n$ 个自由度的完整动力系统的“位形空间” $V_n$，或“相空间” $W_{2n}$。老书的表述虽然不够精确，但很清楚：对许多与 $V_n$ 有关的问题，我们只应考虑定义在开集 $U\subset V_n$ 上的局部函数 $f:U\to\mathbb R$。例如，拉格朗日坐标函数 $q_i$（$i=1,2,\ldots,n$）只在某个坐标卡上局部定义；拉格朗日运动方程
\[
\frac{\mathrm d}{\mathrm dt}\left(\frac{\partial L}{\partial\dot q_i}\right)
-\frac{\partial L}{\partial q_i}=0
\]
也只在某些局部坐标系 $(q_1,\ldots,q_n)$ 中成立。要从整体上考察动力系统的行为，就必须把对应于不同开集 $U$ 的局部函数拼合起来。

为此，先注意到：对于每个 $U$，相应函数集合 $\{f:U\to\mathbb R\mid U\subset V_n\}$ 在逐点加法、乘法下构成含幺交换环，记为 $\Gamma(U)$。若 $V\subset U$，则有自然的限制映射
\[
r_V^U:\Gamma(U)\longrightarrow\Gamma(V),\qquad
r_V^U(\varphi)=\varphi|_V:V\longrightarrow\mathbb R.
\]
换句话说，开集上的局部 $C^\infty$ 可微函数构成一个“预层”（参见第三章）。

接下来，需要考虑函数 $f:U\to\mathbb R$ 在任一点 $x\in U$ 的“芽”。设 $f:U\to\mathbb R$ 和 $g:V\to\mathbb R$ 是局部函数；如果它们在某个含有 $x$ 的开子集 $W\subset U\cap V$ 上相等，就称 $f$ 与 $g$ 在 $x$ 处等价，记作 $f\sim g$。函数 $f$ 在点 $x\in W$ 的芽记作 ${}_xf$，就是此等价关系确定的函数等价类。这个定义事实上已隐含在初等单复变函数论中。

容易看出：在固定点 $x$ 处的这些函数芽构成一个局部环，这里的“局部环”取其现代技术意义。因此，再补上适当的拓扑构造，就可以把“$M$ 上局部 $C^\infty$ 可微函数芽的层” $\mathcal O_M$（正文对拓扑空间 $X$ 记作 $\mathcal O_X$）定义为各点局部环 $\mathcal O_{M,x}$ 的不交并
\[
\coprod_{x\in M}\mathcal O_{M,x}.
\]
于是，经典力学中的微分流形 $V_n$ 或 $W_{2n}$，乃至任何微分流形，

% PDF 10
\sourcepage{6}{10}
都是局部赋环空间 $(X,\mathcal O_X)$ 的例子：也就是配有结构层 $\mathcal O_X$ 的拓扑空间 $X$。

术语虽然显得沉重，芽构成的层其实是一个古老的想法，不过通过一种新而有用的构造得到了精确表述。这些抽象观念并非生来就抽象。虽然可以用抽象方式引入它们，也完全可以通过处理经典例子，以“实验”的方式发现它们。

引入层，是为了在局部性质与整体性质之间建立联系。在这方面，层的引入简化了对经典运动方程解曲线的整体研究，而这本来是一个困难的问题。如果我们同意利希内罗维茨的看法，认为微分流形最具体的模型就是“位形空间”，那么层论似乎仅仅出于技术或审美上的需要，人为地把具体问题抽象化了。然而，按照现代观点，当基域限于 $\mathbb R$ 或 $\mathbb C$ 时，非奇异代数簇可以看作代数几何中的特例。甚至在抽象情形下，塞尔也已在著名论文 FAC 中证明，可以利用扎里斯基拓扑，把层论思想推广到更一般的抽象代数簇。这里塞尔遵循的是嘉当的想法：用 $V$ 上局部正则函数芽的层，取代不可约代数簇 $V$ 的“有理函数域”；此时 $V$ 已不必不可约。格罗滕迪克的概形也是局部赋环空间 $(X,\mathcal O_X)$。\footnote{Demazure--Gabriel 把这些空间称为“几何空间”，因为没有 $\mathcal O_X$ 的空间，似乎还不足以引起几何上的兴趣。}

不过，也不能回避这里出现的新困难：微分流形是豪斯多夫空间，而概形 $S$ 甚至未必是 $T_1$ 空间，只能保证是 $T_0$ 空间。也就是说，对于两个不同的点 $x,y\in S$，其中一点总有一个不含另一点的开邻域，但这种关系对 $x,y$ 并不对称。换句话说，单点 $x\in S$ 未必闭；$\{x\}$ 的闭包 $\overline{\{x\}}$ 可能很大。第二卷还会回到这个问题。本卷主要讨论构成概形的基本部分，即所谓仿射概形（第三章），或仿射代数空间（第一章）。

\Section{1}{预备知识}
我们假定读者熟悉群、环、域等基本代数结构。BOURBAKI-COMMUTATIVE ALGEBRA 各卷包含这里所需的一切。ZARISKI--SAMUEL 的著作虽然有些旧式，也很有用。还假定读者熟悉一般拓扑的基本知识，包括拓扑空间、连续映射，以及层论。层论的标准参考书是 GODEMENT。[关于这一主题的简短介绍，也可参见1973年5月在纽约州立大学布法罗分校讲授、以 SGA-4 为启发的 TOPOI 短期课程。]

此外，读者应熟悉范畴论语言，也就是范畴的定义、从一个范畴到另一个范畴的函子、两个给定范畴之间的函子范畴，以及或许还包括伴随函子的概念。[从课程一开始，从范畴 $C$ 到集合范畴的可表函子，以及协变可表函子 $C\to\Sets$ 组成的范畴，就起着重要作用。我在第一章第1节补入了理解格罗滕迪克讲授所需的最少知识。]

\Section{2}{第一卷提要}
尽管格罗滕迪克带来了种种变革，代数几何仍然是一门“方程的几何理论”。第一章会清楚地说明这一点。我们从一个十分一般的多项式方程组
\[
S=\{f_j(T_i)=0\mid i\in I,\ j\in J\}
\]
出发，其中指标集 $I,J$ 任意，系数属于一个含幺交换基环 $k$。\footnote{这里所考虑的环一律是含幺交换环；任一环同态 $f:A\to B$ 都保持幺元，即 $f(1_A)=1_B$。参见第一章第2节。}我们考虑的解 $(a_i)_{i\in I}$，其坐标 $a_i$ 可以属于

% PDF 11
\sourcepage{7}{11}
任意 $k$-代数（参见第一章第2节）；尤其不把解限制在 $k^I$ 中。这样，每个 $k$-代数 $k'$ 都对应一个子集 $V_S(k')\subset k'^I$，而每个 $k$-代数同态 $f:k'\to k''$ 都对应一个映射
\[
V_S(f):V_S(k')\longrightarrow V_S(k'').
\]
因此，用技术性的语言说，我们得到一个从 $k$-代数范畴 $\G_k$ 到集合范畴的协变函子。\footnote{这个很简单的例子，可以在很早的阶段就用来引入“范畴”或函子语言，因为每个高中毕业生都应知道什么是方程组。至于为什么 $V_S$ 应当是函子而不是集合，第一章会加以讨论。}第一、二章专门研究这个函子。研究有内蕴和外在两种观点，区别在于是否希望忘掉所有仿射嵌入。我们主要关心的，远远是内蕴的观点。因此，两个方程组 $S,S'$，即使其指标对 $(I',J')\ne(I,J)$，也当且仅当相应函子 $V_S,V_{S'}$ 等价时，才被视为等价。

固定 $I$ 后，可以把问题归结到 $S$ 是多项式环 $P_I=k[T_i]_{i\in I}$ 的一个理想 $\mathfrak J$ 的情形。此时，$\mathfrak J\ne\mathfrak J'$ 蕴含 $V_{\mathfrak J}\ne V_{\mathfrak J'}$。但若要避开“周围的仿射空间” $\mathbb E^I$（“$k'^I$”），还需要引入代数 $A=P_I/\mathfrak J$。主要结果是：我们所要找的内蕴函子 $\G_k\to\Sets$，恰好就是协变可表函子
\[
V_A:k'\longmapsto\Hom_{k\text{-alg}}(A,k'),
\]
它给变动的 $k'$ 配上从 $A$ 到 $k'$ 的 $k$-代数同态的集合。\footnote{我们再次鼓励对范畴还不熟悉的读者，在这个非常具体的情形中学习可表函子的概念。这里显然没有用到范畴 $\G_k$ 的任何特殊性质。进一步说明见第一章第1节。}这样的函子记作 $V_A$、$\mathfrak X_A$，或简称 $\mathfrak X$，称为由 $k$-代数 $A$ 表示的 $k$ 上仿射代数空间。\footnote{字母 $V$ 表示（代数）“簇”（variety）。虽然我们不系统讨论经典代数几何这些旧有的研究对象，它们仍然居于相当重要的中心地位。}\footnote{当 $k=\mathbb Z$ 时，$\mathfrak X$ 称为绝对仿射代数空间，也就是 $\mathbb Z$ 上的仿射代数空间。}具体说来：$\mathfrak X_A$ 的一个“内蕴的” $k'$-值点 $P$，就是一个 $k$-代数同态 $P:A\to k'$，即

% PDF 12
\sourcepage{8}{12}
\[
\begin{gathered}
P(a+b)=P(a)+P(b),\qquad P(ab)=P(a)P(b),\\
P(1)=1,\qquad P(\lambda a)=\lambda P(a)\quad(\lambda\in k).
\end{gathered}
\]
\footnote{如果觉得这“太抽象”，就先考虑一个嵌入 $\mathfrak X_A(k')\hookrightarrow k'^I$，核对点 $P$ 的坐标就是它在一组生成元 $\xi_i$ 上的取值 $P(\xi_i)$，然后再忘掉坐标。}
由此得到三个显然的事实：
\begin{enumerate}[label=\alph*)]
\item $k'$-值点的集合 $\mathfrak X_A(k')$ 就是 $\Hom_{k\text{-alg}}(A,k')$，即
\[
\mathfrak X_A(k')=\Hom_{k\text{-alg}}(A,k').
\]
\item 对应 $k'\mapsto\mathfrak X_A(k')$ 定义了协变函子 $\G_k\to\Sets$。\footnote{译注：原书这里将定义域误排为 $\G_A$，依上下文订为 $\G_k$。}
\item 每个 $k$-代数同态 $h:A\to B$ 都对应一个反方向的映射
\[
h^*:\mathfrak X_B(\ )\longrightarrow\mathfrak X_A(\ ),
\]
由唯一合理的复合方式定义。
\end{enumerate}
a)、b) 以内蕴方式重新定义了解函子 $k'\mapsto\mathfrak X_A(k')$：$\mathfrak X_A$ 是由 $A$ 表示的（协变）可表函子 $\G_k\to\Sets$。c) 则说明：协变可表函子 $\mathfrak X_A$（$A\in\Ob\G_k$）构成的范畴，与 $\G_k$ 的反范畴 $\G_k^{\op}$ 等价；对应 $A\mapsto\mathfrak X_A$ 定义了一个反变函子。\footnote{这个性质对任何范畴 $C$ 都成立。}

以上这些范畴性质，并没有告诉我们关于 $k$-代数范畴及其“几何解释”的任何特殊信息。\footnote{参见上一脚注。}与之相对，在第三章中，我们将给函子 $\mathfrak X_A$ 配上一个几何对象，也就是由 $A$ 定义的仿射概形 $(X,\mathcal O_X)$。它是所谓局部赋环空间的一种特例。\footnote{赋环空间 $(X,\mathcal O_X)$ 是由拓扑空间 $X$ 及其上的结构层 $\mathcal O_X$ 组成的二元组。现代数学中的大多数几何结构都是赋环空间，例如各种流形（拓扑流形、微分流形、解析流形）：它们都是配有相应类型（连续、可微、解析等）局部函数芽层的拓扑空间 $X$。“局部赋环”是指各个茎都是局部环。详见第三章第7节。}在概述这件事之前，

% PDF 13
\sourcepage{9}{13}
我先极简要地回顾一下 $k$-代数的几何意义，以便切实建立旧式代数几何、整个几何学和当代应用之间的联系。

在经典代数几何中，$k$ 是复数域，即 $k=\mathbb C$，我们在 $\mathbb C^n$ 中找解，其中 $n=\#I$。\footnote{指标集 $J$ 是有限还是无限并不重要：由希尔伯特基定理，$S$ 生成的理想 $\mathfrak J$ 总有有限基。}集合 $V=V_S(\mathbb C)$ 被定义为一个“代数簇”。多项式函数 $f:\mathbb C^n\to\mathbb C$ 在 $V$ 上的限制 $f|_V$，组成一个有限生成的 $\mathbb C$-代数 $A$。由希尔伯特零点定理，
\[
A\simeq\mathbb C[T_1,\ldots,T_n]/\sqrt{\mathfrak J},
\]
其中 $\sqrt{\mathfrak J}$ 表示理想 $\mathfrak J$ 的根。\footnote{环 $A$ 的理想 $\mathfrak a$ 的根是 $\sqrt{\mathfrak a}=\{a\in A\mid \exists m\in\mathbb Z_{>0},\ a^m\in\mathfrak a\}$；显然 $\sqrt{\mathfrak a}\supset\mathfrak a$。当且仅当 $\mathfrak a=\sqrt{\mathfrak a}$ 时，称 $\mathfrak a$ 为根理想。特别地，$\sqrt{(0)}$ 是 $A$ 的幂零根，本课程记作 $\operatorname{Nil}(A)$；$A$ 约化，当且仅当 $\operatorname{Nil}(A)=0$，也就是 $A$ 中唯一的幂零元是零。}\footnote{零点定理的原始表述，对任意代数闭域 $k$ 都成立：若 $f$ 在理想 $\mathfrak a$ 的所有“零点”（Nullstellen）上消失，则存在正整数 $m$ 使 $f^m\in\mathfrak a$。}所以 $A$ 没有非零幂零元。

反过来，若 $A$ 是有限生成的约化 $\mathbb C$-代数，就可以把 $A$ 与某个代数簇上的复值多项式函数代数等同起来。例如，在 $\mathbb C^3$ 中，若 $V\ne\mathbb C^3$，则 $V$ 由有限多个不可约曲面、有限多条不可约曲线和有限多个点组成。特别地，$V$ 不可约，当且仅当 $A$ 是整环。\footnote{含幺交换环 $A$ 是整环，当且仅当 $A\setminus\{0\}$ 是乘法闭的，也就是 $1\in A\setminus\{0\}$，且 $a,b\ne0$ 蕴含 $ab\ne0$。}
\begin{center}
\begin{tikzpicture}[scale=.8]
\draw[->] (0,0)--(-.8,-.7) node[below]{$x$};
\draw[->] (0,0)--(4.4,0) node[right]{$y$};
\draw[->] (0,0)--(0,3.1) node[above]{$z$};
\node[below right] at (0,0){$0$};
\draw[thick] (.4,1.2) .. controls (.7,.5) and (1.7,.7) .. (2.2,1.6) .. controls (2.7,2.5) and (3,2.9) .. (3.7,2.9);
\fill (.7,2.5) circle (1.6pt); \fill (1,1.9) circle (1.6pt); \fill (1.8,2.4) circle (1.6pt);
\draw (2.8,.7) .. controls (3.2,1.2) and (3.3,1.9) .. (3.1,2.2);
\draw (3.9,.6) .. controls (3.7,1.2) and (3.9,2) .. (4.1,2.3);
\draw (3.15,1.4) arc[start angle=180,end angle=360,x radius=.35,y radius=.12];
\draw[dashed] (3.15,1.4) arc[start angle=180,end angle=0,x radius=.35,y radius=.12];
\end{tikzpicture}
\end{center}
在经典代数几何中，带有非零幂零元的有限生成 $\mathbb C$-代数，完全没有几何地位。\footnote{环 $A$ 的元素 $f$ 称为幂零元，是指某个正整数 $m$ 满足 $f^m=0$。应用数学者、工程师等每次略去 $h$ 阶及以上的“无穷小量”，写下 $(x-a)^h=0$ 之类的式子时，实际上都引入了幂零元。}格罗滕迪克反对这种看法，因为这样的

% PDF 14
\sourcepage{10}{14}
代数表示的正是无穷小对象（参见第一章第13节）。

因此，为培养几何直觉，读者应把任意 $k$-代数理解为 $V$ 上多项式函数代数的自然推广：这里基环 $k$ 任意，代数也不一定有限生成。问题在于“如何从 $A$ 以某种方式恢复 $V$”。

若 $W\subset V$，则多项式函数的限制给出相应多项式函数环之间的满同态 $A_V\to A_W$。然而，既然非约化代数没有地位，经典代数几何就不可能是“函子式”的。那时虽没有明确地定义函子，却已在暗中“使用”它们，只不过必须反复写出往往长得令人难以忍受的陈述。

给任意 $A\in\Ob\G_k$ 所配的几何对象 $(X,\mathcal O_X)$，也就是相应的仿射概形，或者等价地，由 $A$ 表示的仿射代数空间 $\mathfrak X_A$，是对代数簇这一旧概念的极有力的精细化。

空间 $X$ 是 $A$ 的谱 $\Spec A$，即 $A$ 的全体素理想组成的集合；$X=\varnothing$ 当且仅当 $A=0$。$X$ 的拓扑由 $A$ 的根理想来定义，其方式受到仿射空间的扎里斯基拓扑启发，参见第一章第14节。仅凭 $\Spec A$ 还不能恢复 $A$，因为 $\Spec A$ 与 $\Spec(A/\operatorname{Nil}A)$ 同胚。因此，若只是构造谱，局限于约化代数便已足够。结构层 $\mathcal O_X$ 与 $X$ 合在一起，则使我们能够恢复 $A$，因为
\[
A\simeq\Gamma(X,\mathcal O_X)=H^0(X,\mathcal O_X).
\]
也就是说，$A$ 同构于层 $\mathcal O_X$ 的整体截面环；或者说，$A$ 就是 $X$ 以 $\mathcal O_X$ 为系数的零次上同调 $k$-代数。结构层 $\mathcal O_X$ 的

% PDF 15
\sourcepage{11}{15}
定义相当技术性，不过这里可以指出：它在每个点 $\mathfrak p\in\Spec A$ 的茎，就是局部环 $A_{\mathfrak p}$。事实上，一切构造在含幺交换环范畴 $\G$ 中都有意义。注意，$k$ 在构造中并不起作用。\footnote{这在心理上似乎有些令人不安：原来方程组的系数环竟然消失了！不过并没有完全消失：总有同态 $h:\mathbb Z\to A$，$n\mapsto n1_A$，所以 $A$ 含有一个环 $\mathbb Z/m\mathbb Z$，其中 $(m)=\Ker h$，$m$ 是 $A$ 的特征。即便反对格罗滕迪克代数几何方法的人，也承认概形对算术问题特别有用。$\mathbb Z$ 起着普遍性的作用。}

\BookChapter{1}{多项式方程组解集的函子描述}
\begin{quote}
问：我们理解您对专家知识的忧虑。不过，说起来，如果要向外行解释什么是代数几何，我觉得恩里克斯那本老书的书名至今仍然合适。\footnote{恩里克斯：“方程的几何理论……”；理解本课程并不需要读这本书。}您怎么看？

格罗滕迪克答：是的！但您的“外行”应该知道什么是代数方程组。柏拉图要学懂这个，可得花上好几年！

问：两千年都过去了，我们是不是也该有点信心，相信每个优秀的高中毕业生都能理解什么是仿射概形呢？您怎么看？……？？
\begin{flushright}摘自一次与格罗滕迪克有关“生存”的简短谈话。\end{flushright}
\end{quote}
\textbf{提要。}设 $S=\{f_j(T_i)\}$（$i\in I,j\in J$）表示任意一个多项式方程组 $f_j(T_i)=0$，\footnote{参见全书提要。}其系数属于含幺交换基环 $k$。我们研究它在任意 $k$-代数 $k'$ 中的解集 $V_S(k')$，其中 $k'\in\Ob\G_k$。\footnote{参见第2节；$\G_k$ 表示 $k$-代数范畴。当 $k=\mathbb Z$ 为整数环时，$\G_{\mathbb Z}=\G$ 是含幺交换环范畴。}对应 $k'\mapsto V_S(k')$ 定义一个协变函子 $\G_k\to\Sets$。主要问题是：不依赖于任何仿射嵌入 $V_S(k')\hookrightarrow\mathbb E^I$，在等价意义下刻画这些函子。答案是：函子 $\G_k\to\Sets$ 等价于某个 $V_S$，当且仅当它等价于某个函子 $\mathfrak X_A$，后者定义为
\[
k'\longmapsto\mathfrak X_A(k')=\Hom_{k\text{-alg}}(A,k')
\qquad(k'\in\Ob\G_k),
\]
% PDF 16
\sourcepage{12}{16}
其中 $A=P_I/\mathfrak J$，$\mathfrak J=P_I\cdot S$ 是多项式环 $P_I=k[T_i]_{i\in I}$ 的一个理想。我们说 $\mathfrak X_A$ 由 $A$ 表示；于是，函子 $\G_k\to\Sets$ 等价于某个 $V_S$，当且仅当它可表，参见第1节。\footnote{第1节包含有关可表函子的一切必要预备知识，取自格罗滕迪克在布法罗讲授的拓扑斯课程和代数群课程。}$\mathfrak X_A$ 称为 $k$ 上的仿射代数空间。特别地，当 $A=P_I$ 时，$\mathfrak X_{P_I}=\mathbb E^I$，即 $I$ 型标准仿射空间。函子 $\mathfrak X_A$ 构成范畴 $\Aff_k$ 的对象，而这个范畴与 $\G_k$ 的反范畴 $\G_k^{\op}$ 等价；这是一个纯粹的范畴事实。特别地，当 $k=\mathbb Z$ 时，$\Aff_{\mathbb Z}=\Aff$ 是“绝对”（即 $\mathbb Z$ 上）仿射代数空间范畴，与含幺交换环范畴 $\G$ 相反。

第一章分为两部分。第一部分证明 $\Aff_k\simeq\G_k^{\op}$；第二部分讨论 $\G$ 的某些特殊子范畴，它们具有重要的几何与历史意义。

\textbf{注。}第一章紧随录音展开，内容几乎与 Springer 新版 EGA-I 的引言相同；EGA-Springer 的后文并不再使用这些材料。\footnote{马宁的讲义也从同样的方法入手，不过所给的细节比本课程少得多。}第1、2节补入了可表函子和 $k$-代数的回顾；格罗滕迪克的讲授是从第3节开始的。

\PartTitle{第一部分}{同构 $\Aff_k\simeq\G_k^{\op}$}
\textbf{提要。}为了把 $V_S$ 刻画为可表函子，我们先建立等价 $V_S\simeq V_{\mathfrak J}$，其中 $\mathfrak J=P_I\cdot S$，参见第6节；再利用 $A$ 的泛性质，证明 $V_{\mathfrak J}\simeq\mathfrak X_A$，其中 $A=P_I/\mathfrak J$，参见第7节。反过来，任意 $A\in\Ob\G_k$ 都可表成这样的商，因此每个 $\mathfrak X_A$ 都来自某个 $V_{\mathfrak J}$。

\Section{1}{可表函子、函子范畴与相对范畴}
\textbf{函子范畴 $\underline{\Hom}(C,\Sets)$、$\underline{\Hom}(C^{\op},\Sets)$，以及范畴 $S/C$、$C/S$。}\footnote{本节主要取自布法罗的 AAG 课程。出于几个原因，尤其是为避免不同假设混杂所造成的混乱，我选择自己重写整个主题。进一步阅读可参见 D-ADVANCES, II，第378页（不含证明），或 FONDEMENTS，第195-01页。正式论述见 EGA-Springer，第0章第1节，第19页；原始文献是 TOHOKU 和 SGA-3, I。}
设 $A\in\Ob C$ 是范畴 $C$ 的一个固定对象。对应
\[
X\longmapsto\Hom_C(A,X)\tag{1.1}
\]

% PDF 17
\sourcepage{13}{17}
给每个 $X$ 配上从 $A$ 到 $X$ 的 $C$-态射集合，\footnote{对于任意一对对象 $X,Y\in\Ob C$，$\Hom_C(X,Y)$ 都是一个集合。这里不讨论范畴论的基础问题。有些作者把 $\Ob C$ 看作一个类，不一定是集合；或者遵循格罗滕迪克在 SGA-3 中的方法，把它放在一个满足特定性质、称为“宇宙”的大集合中。}它对于 $X$ 是函子性的，也就是定义了协变函子 $V_A:C\to\Sets$，称为由 $A$ 表示的协变函子 $C\to\Sets$。

事实上，式（1.1）定义了对象之间的映射 $\Ob C\to\Ob\Sets$。此外，若 $f:X\to Y$ 是 $C$ 中的态射，就有自然映射
\[
\Hom_C(A,X)\xrightarrow{\ V_A(f)\ }\Hom_C(A,Y),\tag{1.2}
\]
由与 $f$ 左复合来定义，即 $u\mapsto f\circ u$，其中 $u\in\Hom_C(A,X)$。

类似地，可以定义反变函子 $h_A:C^{\op}\to\Sets$：\footnote{$C^{\op}$ 是 $C$ 的反范畴；它具有与 $C$ 相同的对象和箭头，不过箭头方向以及复合的次序均反转。}
\[
h_A(X)=\Hom_C(X,A).\tag{1.3}
\]
对每个 $C$-态射 $f:X\to Y$，映射 $h_A(f):h_A(Y)\to h_A(X)$ 定义为
\[
h_A(f)(v)=v\circ f\qquad\bigl(v\in\Hom_C(Y,A)\bigr).\tag{1.4}
\]
于是 $h_A:C^{\op}\to\Sets$ 就是由 $A$ 表示的从 $C$ 到集合的反变函子。

\textbf{注。}
\begin{enumerate}[label=\arabic*)]
\item 引入反范畴 $C^{\op}$ 后，$V_A$ 与 $h_A$ 两种情形当然可以彼此化归。
\item 两类函子 $V_A,h_A$ 都会经常而自然地出现。本课程所需的主要例子来自代数几何，此时
\end{enumerate}

% PDF 18
\sourcepage{14}{18}
$C$ 是 $k$-代数范畴 $\G_k$，其中 $k$ 为固定的含幺交换基环，参见第2节；或者 $C$ 是它的对偶范畴 $\G_k^{\op}$。\footnote{把 $\G_k^{\op}$ 自然地与所谓 $k$ 上的仿射代数空间范畴等同起来，正是本章的主要结果。}事实上，元素（环同态）
\[
u\in\Hom_{\G_k}(A,k')\qquad(k'\in\Ob\G_k)
\]
就是由 $A$ 表示的仿射代数空间的 $k'$-值点。\footnote{直观地说，就是“坐标属于 $k'$ 的点”；不过，我们想要的是摆脱坐标。}

从 $C$ 到 $\Sets$ 的协变（反变）函子 $F$，称为由对象 $A\in\Ob C$ 表示，是指 $F$ 等价于某个 $V_A$（或 $h_A$）。若是如此，表示对象 $A$ 在同构意义下唯一确定。为明确说明这一陈述，需要引入从 $C$ 到 $\Sets$ 的协变（反变）函子范畴 $\underline{\Hom}(C,\Sets)$（或 $\underline{\Hom}(C^{\op},\Sets)$），把上面的定义形式化。

按照熟知的构造，\footnote{若 $\mathcal G,\mathcal B$ 是两个范畴，则 $\operatorname{Funct}(\mathcal G,\mathcal B)=\underline{\Hom}(\mathcal G,\mathcal B)$ 是协变函子 $\mathcal G\to\mathcal B$ 的范畴，以自然变换为态射。}$\underline{\Hom}(C,\Sets)$ 的对象是协变函子 $F:C\to\Sets$。若 $F,G\in\Ob\underline{\Hom}(C,\Sets)$，则态射 $u:F\to G$ 是一个自然变换，即：对每个 $A\in\Ob C$，都有映射
\[
u(A):F(A)\longrightarrow G(A),
\]
并且对每个 $C$-态射 $f:A\to B$，下图交换：
\[
\begin{tikzcd}[ampersand replacement=\&]
F(A)\arrow[r,"F(f)"]\arrow[d,"u(A)"']\&F(B)\arrow[d,"u(B)"]\\
G(A)\arrow[r,"G(f)"']\&G(B).
\end{tikzcd}\tag{1.5}
\]
特别地，$u$ 是等价，当且仅当对任意 $A\in\Ob C$，竖直箭头 $u(A)$ 都是集合之间的等价，也就是双射。

% PDF 19
\sourcepage{15}{19}
$\underline{\Hom}(C^{\op},\Sets)$，即取值于集合的反变函子范畴，其情形留给读者处理。

把 $\underline{\Hom}(C^{\op},\Sets)$ 记为 $\widehat C$；相应地，对反范畴应用同一构造，得到 $\widehat{C^{\op}}=\underline{\Hom}(C,\Sets)$。\footnote{译注：原书以帽号及上标圆圈表示这两个函子范畴；这里将上标明确写在帽号内，以区别于函子范畴本身的反范畴。}$\widehat C$ 是 $C$ 的一个自然扩充：它把每个对象 $X\in\Ob C$ 与“由 $X$ 表示”的反变函子 $h_X$ 等同起来。这个函子从 $C$ 到集合，定义为 $h_X(Y)=\Hom_C(Y,X)$。事实上，由
\[
i:C\hookrightarrow\widehat C\tag{1.6}
\]
及
\[
i(X)=h_X\qquad(X\in\Ob C)\tag{1.7}
\]
定义的嵌入函子，是从 $C$ 到 $\widehat C$ 的全忠实函子，\footnote{函子 $F:C\to\mathcal D$ 称为忠实函子，是指对每一对 $A,B\in\Ob C$，所诱导的映射 $\Hom(A,B)\to\Hom(F(A),F(B))$ 都是单射。若这些映射对任意 $A,B$ 都是双射，则称为全忠实函子。原书还把 $F(C)$ 视作 $\mathcal D$ 的子范畴；在全忠实时，它是满子范畴。}因而可以把 $C$ 看作 $\widehat C$ 的满子范畴。\footnote{设 $\mathcal B$ 是范畴 $\mathcal G$ 的子范畴。对 $X,Y\in\Ob\mathcal B$，有自然单射 $\Hom_{\mathcal B}(X,Y)\hookrightarrow\Hom_{\mathcal G}(X,Y)$，即 $\mathcal B$ 中的每个态射也是 $\mathcal G$ 中的态射。当且仅当对每一对这样的对象都有 $\Hom_{\mathcal B}(X,Y)=\Hom_{\mathcal G}(X,Y)$ 时，称 $\mathcal B$ 为 $\mathcal G$ 的满子范畴。}$C$ 在 $i$ 下的本质像中的对象，\footnote{若 $i:\mathcal E\to\mathcal F$ 全忠实，则 $\mathcal E$ 在 $\mathcal F$ 中的本质像，是 $\mathcal F$ 中与 $i(\mathcal E)$ 的对象同构的那些对象所组成的满子范畴。也就是说，$F\in\Ob\mathcal F$ 属于本质像，当且仅当存在 $X\in\Ob\mathcal E$ 使 $F\simeq i(X)$。}就是所谓反变可表函子。

本章大量使用的协变可表函子，可以通过对反范畴 $C^{\op}$ 应用帽号构造来获得。换句话说，可以定义一个嵌入

% PDF 20
\sourcepage{16}{20}
\[
j:C^{\op}\hookrightarrow\widehat{C^{\op}},\tag{1.8}
\]
把 $C^{\op}$ 嵌入从 $C$ 到集合的协变函子范畴：令 $j(X)=j_X$，其中 $j_X$ 定义为 $Y\mapsto j_X(Y)=\Hom_C(X,Y)$。从 $C$ 到 $\Sets$ 的协变可表函子，恰好是 $C^{\op}$ 在 $j$ 下本质像中的对象。

以上讨论可以归纳如下：从 $C$ 到集合有两种函子 $h_X=\Hom_C(\ ,X)$ 和 $j_X=\Hom_C(X,\ )$，其对象映射为
\[
\begin{aligned}
Y&\longmapsto h_X(Y)=\Hom_C(Y,X),\\
Y&\longmapsto j_X(Y)=\Hom_C(X,Y),
\end{aligned}\tag{1.9}
\]
并以唯一可能的复合方式处理 $C$-态射 $f:Y\to Z$。因此，$h_X$ 是反变函子 $C^{\op}\to\Sets$，$j_X$ 是协变函子 $C\to\Sets$。分别称它们为由 $X$ 表示的从 $C$ 到集合的反变、协变函子。

更一般地，从 $C$ 到 $\Sets$ 的反变（协变）函子 $F$ 称为可表函子，是指存在对象 $X\in\Ob C$，使 $F\simeq h_X$（或 $F\simeq j_X$）。无论哪种情形，都说 $X$ 是范畴 $C$ 中表示函子 $F$ 的对象。显然，若 $X'\simeq X$，则 $X'$ 表示 $F$ 当且仅当 $X$ 表示 $F$。换句话说，表示对象 $X$ 在同构意义下确定。\footnote{这个简单的核实留给读者。}

\textbf{注。}多数时候，直接把由 $X$ 表示的反变（协变）函子 $F$ 与 $h_X$（或 $j_X$）等同起来，上面的概念就已足够。\footnote{只要不要求同构的唯一性，这样做并不引起问题。}

% PDF 21
\sourcepage{17}{21}
在更精细的问题中，则有必要强调某个特定元素 $\xi\in F(X)$ 的选择：它是恒等态射
\[
1_X\in\Hom_C(X,X)=h_X(X)=j_X(X)
\]
在相应集合等价 $\gamma_X(X)$ 下的像。事实上，在反变情形中，$C$ 中的态射 $f:Y\to Z$ 诱导交换方块
\[
\begin{tikzcd}[ampersand replacement=\&]
h_X(Z)\arrow[r,"h_X(f)"]\arrow[d,"\gamma_X(Z)"']\&h_X(Y)\arrow[d,"\gamma_X(Y)"]\\
F(Z)\arrow[r,"F(f)"']\&F(Y)
\end{tikzcd}
\qquad
\begin{tikzcd}[ampersand replacement=\&]
1_X\arrow[d,mapsto,"\gamma_X(X)"]\\
\xi
\end{tikzcd}\tag{1.10}
\]
特别地，若 $Z=X$，就有映射 $h_X(f):h_X(X)\to h_X(Y)$，并满足 $h_X(f)(1_X)=f$。元素 $\xi$ 具有如下泛性质：
\[
\gamma_X(Y)(f)=F(f)(\xi).\tag{1.11}
\]
因此，当我们用具体的等价 $Y\mapsto\gamma_X(Y)$，即
\[
\gamma_X(Y):h_X(Y)\longrightarrow F(Y),\tag{1.11}
\]
来区别 $F$ 与同它等价的 $h_X$ 时，就可以说：$F$ 完全由表示对象 $X\in\Ob C$ 及特定元素 $\xi=\gamma_X(X)(1_X)$ 确定。\footnote{译注：本页两式在原书中均编号为（1.11），这里保留原编号。}

协变情形中对 $\xi$ 的讨论，以及相应的结论——$F$ 由 $X$ 和泛元素 $\xi\in F(X)$ 唯一确定——留给读者。

从这个更具体的角度来看，较准确的定义应当是：

% PDF 22
\sourcepage{18}{22}
函子 $F$ 可表，当且仅当存在一对 $(X,\xi)$，其中 $X\in\Ob C$、$\xi\in F(X)$，而且 $\xi$ 满足上述泛性质。

在前面几段中，对象 $X$ 是固定的。现在让 $X$ 在 $\Ob C$ 中变动，并把上面的可表函子看作两个新范畴的对象；这些范畴的态射是自然变换。格罗滕迪克对两种变性下的这种变化，统称为“换基”。分别讨论如下。

\textbf{反变情形：}
\[
Y\longmapsto\Hom_C(Y,X)=h_X(Y)\qquad(Y\in\Ob C).
\]
随着 $X\in\Ob C$ 变动，这些反变可表函子构成的范畴，自然地与 $C$ 同构，典范对应为 $X\mapsto h_X$。这是一个全忠实函子，因而可以把 $C$ 与反变函子范畴 $\widehat C=\underline{\Hom}(C^{\op},\Sets)$ 的一个满子范畴等同起来。\footnote{在 TOHOKU 中，格罗滕迪克利用这一典范嵌入 $C\hookrightarrow\widehat C$，把逆极限的构造归结为集合情形，参见第二章第1、2节。}

\textbf{协变情形：}
\[
X\longmapsto\Hom_C(S,X).
\]
类似地，可把协变可表函子 $j_X$ 与作为反范畴 $C^{\op}$ 的对象的 $X$ 等同起来。\footnote{可以用 $j$ 来定义归纳极限构造，参见第二章第3、8节。}换句话说，$C$ 的反范畴 $C^{\op}$ 与协变可表函子范畴典范同构。

设 $S$ 是范畴 $C$ 的固定对象。\footnote{译注：原书此处写作固定对象 $X$，下文却统一使用 $S$，这里将其统一为 $S$。}定义两个范畴 $S/C$ 和 $C/S$：它们的对象分别是以 $S$ 为源和以 $S$ 为靶的 $C$-态射，态射则是使相应三角形交换的 $C$-箭头：

% PDF 23
\sourcepage{19}{23}
\[
\begin{aligned}
u\in\Ob(S/C)&\quad\Longleftrightarrow\quad u:S\longrightarrow X_u,\\
u\in\Ob(C/S)&\quad\Longleftrightarrow\quad u:X_u\longrightarrow S.
\end{aligned}
\]
具体而言，$S/C$ 中从 $u$ 到 $v$ 的态射 $\varphi$，以及 $C/S$ 中从 $u$ 到 $v$ 的态射 $\psi$，分别由下列交换图给出：
\[
\begin{tikzcd}[ampersand replacement=\&]
S\arrow[r,"u"]\arrow[dr,"v"']\&X_u\arrow[d,"\varphi"]\\
\&X_v
\end{tikzcd}
\qquad
\begin{tikzcd}[ampersand replacement=\&]
X_u\arrow[r,"u"]\arrow[d,"\psi"']\&S\\
X_v\arrow[ur,"v"']\&
\end{tikzcd}
\]
即 $v=\varphi\circ u$ 和 $u=v\circ\psi$。\footnote{译注：原书两个三角形下的复合式与其箭头标记不一致，这里按图中态射方向写出。}
这里照例以 $\Ob C$、$\operatorname{Arr}C$ 表示 $C$ 的对象类和箭头类。

$C$ 的对象 $\varnothing_C$ 称为 $C$ 的始对象，是指对每个 $X\in\Ob C$，恰好存在一个态射 $\varnothing_C\to X$。显然，$C$ 的任意两个始对象都同构。在集合范畴中，始对象是空集。

对偶地，$C$ 的对象 $e$ 称为终对象，是指对每个 $X\in\Ob C$，恰好存在一个态射 $X\to e$。集合范畴中的单点集是终对象。任意两个终对象都同构。当然，$C$ 的始对象（终对象）在 $C^{\op}$ 中成为终对象（始对象）。

若 $C$ 有始对象 $\varnothing_C$（终对象 $e$），就可用显然的方式把 $C$ 与 $\varnothing_C/C$（或 $C/e$）等同起来。

\Section{2}{$k$-代数范畴 $\G_k$}
我们的出发点，是研究系数属于含幺交换基环 $k$ 的十分一般的多项式方程组的解。但不能把解限制在 $k$ 中，而要在任意 $k$-代数 $k'$ 中找解。下面从范畴观点回顾 $k$-代数的含义。

设 $\G$ 是含幺交换环范畴；$\G$ 中的态射（或“箭头”）是保持幺元的环同态 $f:A\to B$，其中 $A,B\in\Ob\G$。也就是说，$f$ 满足下面的条件（2.1）、（2.2）、（2.3）：

% PDF 24
\sourcepage{20}{24}
\[
f(x+y)=f(x)+f(y),\tag{2.1}
\]
\[
f(xy)=f(x)f(y),\tag{2.2}
\]
\[
f(1_A)=1_B.\tag{2.3}
\]
整数环 $\mathbb Z$ 是 $\G$ 的始对象，参见第1节。事实上，任意箭头 $\varphi:\mathbb Z\to A$ 都把 $\pm(1+\cdots+1)$ 映到 $\pm(m1_A)$，这里有 $m$ 个 $1$；因此 $\varphi$ 唯一，并由条件（2.3）确定。$\varphi$ 未必是单射。其像 $\varphi(\mathbb Z)\simeq\mathbb Z/m\mathbb Z$，其中 $m\ge0$ 是 $A$ 的特征。如果 $A$ 是域，则 $m=0$ 或 $m$ 为素数。

设 $k$ 是 $\G$ 中的任意一个环。\footnote{“环”一律指 $\G$ 的对象，也就是含幺交换环。}在 $\G$ 中以态射 $h:k\to A$ 为对象所构成的范畴，其对象通常简称为 $k$-代数。也就是说，按照一般范畴构造，$k$-代数范畴 $\G_k$ 定义如下：
\begin{enumerate}[label=\arabic*)]
\item $\G_k$ 的一个对象是一个态射 $h:k\to A$。
\end{enumerate}
通常稍宽松地把 $A$ 称为 $k$-代数，把 $h$ 称为结构态射；应始终记住，$h$ 必须是已知的。\footnote{译注：原书本句误将结构态射写成 $k$，这里按前后文写作 $h$。}
\begin{enumerate}[label=\arabic*),start=2]
\item 在范畴 $\G_k$ 中，从 $u:k\to A$ 到 $v:k\to B$ 的态射，是一个使下列三角形交换的态射 $\varphi:A\to B$：
\end{enumerate}
\[
\begin{tikzcd}[ampersand replacement=\&]
\&k\arrow[dl,"u"']\arrow[dr,"v"]\&\\
A\arrow[rr,"\varphi"']\&\&B.
\end{tikzcd}\tag{2.4}
\]
如上所说，态射 $u,v$ 往往是“明确的”，因而以下说法使用得很广泛：对应于 $h:k\to A$，有外乘法 $k\times A\to A$，定义为

% PDF 25
\sourcepage{21}{25}
\[
(\lambda,a)\longmapsto h(\lambda)a=ah(\lambda).\tag{2.5}
\]
特别地，由于 $h(1_k)=1_A$，有 $(1,a)\mapsto a$。这一数乘 $k\times A\to A$ 作为附加结构，与 $(+,\cdot,1_A)$ 配在一起，并满足
\[
\lambda(a+b)=\lambda a+\lambda b,\tag{2.6}
\]
\[
\lambda(ab)=(\lambda a)b=a(\lambda b).\tag{2.7}
\]
按照这种说法，$k$-代数范畴 $\G_k$ 的对象是 $k$-代数，也就是配有结构态射 $h:k\to A$ 的环 $A\in\Ob\G$；而 $k$-代数态射 $\varphi:A\to B$，是与数乘交换的 $\G$-态射：
\[
\varphi(\lambda a)=\lambda\varphi(a)
\qquad(a\in A,\ \lambda\in k).\tag{2.8}
\]

对于给定的 $k$，$k$-代数是环的特殊情形，乍看起来，研究 $k$-代数似乎会“限制更多”。其实不然：$\mathbb Z$ 是 $\G$ 的始对象，而且每个保持幺元的环同态 $A\to B$ 都与整数数乘交换。换句话说：

\textbf{含幺交换环及保持幺元环同态所构成的范畴 $\G$，与 $\mathbb Z$-代数及 $\mathbb Z$-代数同态所构成的范畴 $\G_{\mathbb Z}$ 完全相同。}

\Section{3}{将 $k'^I$ 中的点与 $k$-代数同态等同}
\textbf{提要。}若 $k'\in\Ob\G_k$，参见第1节，且 $I$ 是任意指标集，则 $k'$ 上 $I$ 型标准仿射空间中的每一点 $x\in k'^I$，都确定多项式环 $P=P_I=k[T_i]_{i\in I}$ 的一个 $k$-代数同态 $\varepsilon_x:P_I\to k'$，其中 $T_i$ 是以 $I$ 为指标、添加到 $k$ 上的

% PDF 26
\sourcepage{22}{26}
一族代数独立的未定元。反过来，任意这样的同态 $u$ 都确定唯一的点 $(u(T_i))_{i\in I}\in k'^I$。

先从代数几何或算术中的通常情形说起。设 $k$ 是含幺交换环。经典情形中的 $k$ 是实数域 $\mathbb R$ 或复数域 $\mathbb C$；但在许多算术问题中，$k$ 可以是整数环 $\mathbb Z$，或者某个数域的整数环。无论如何，暂且固定 $k$，并称之为基环。

设 $P=P_I$ 是以 $k$ 为系数、以集合 $I$ 为指标的未定元族 $T_i$ 中的多项式环：\footnote{不要求 $I$ 一定有限。许多经典情形确实是 $I=\{1\}$、$\{1,2\}$、$\{1,2,3\}$ 等，分别对应仿射直线、平面、三维空间中的方程组。后文会从技术上说明采用这一一般性的理由。}
\[
P=P_I=k[T_i]_{i\in I}.\tag{3.1}
\]
显然，$P_I\in\Ob\G_k$，即 $P_I$ 是由 $T_i$（$i\in I$）生成的 $k$-代数。

设 $x=(x_i)_{i\in I}$ 是一族以 $I$ 为指标的 $k$ 中元素，即每个 $x_i\in k$。称 $x$ 为 $I$ 型标准仿射空间 $k^I$ 的一个点。\footnote{$k^I$ 是所有映射 $I\to k$ 组成的集合；当 $I=\{1,2,\ldots,n\}$ 时，简记为 $k^n$。用格罗滕迪克的话说，$\mathbb E^I$ 是遗忘函子 $\mathbb E^1$ 的 $I$ 次笛卡儿幂；$\mathbb E^1$ 把每个 $k'\in\Ob\G_k$ 对应到其底集。原书将这里的函数集合称作 power set，意指此笛卡儿幂。}

多项式 $f\in P_I$ 定义一个函数 $k^I\to k$。这个函数有时也用相同字母表示，尽管它未必反过来唯一确定多项式 $f$。为避免这种记号滥用，这里把函数记作 $\widetilde f$：给每个 $x\in k^I$，把 $f$ 中的每个未定元 $T_i$ 换成 $x_i$，所得的 $k$ 中元素 $f(x)$，就是 $\widetilde f(x)$。

更一般地，还可对任意 $k$-代数 $k'$，定义函数 $\widetilde f_{k'}:k'^I\to k'$。

% PDF 27
\sourcepage{23}{27}
固定 $x\in k'^I$。对应 $i\mapsto x_i\in k'$ 通过代入
\[
f(T_i)\longmapsto f(x_i)
\]
定义一个 $k$-代数同态 $u:P_I\to k'$，参见第2节式（2.8）。这也可写作 $x\mapsto\varepsilon_x$，其中 $x\in k'^I$、$\varepsilon_x\in\Hom_{k\text{-alg}}(P_I,k')$。多项式 $f\in P_I$ 在这个同态下的像 $\varepsilon_x(f)$，就是把 $f$ 中的每个 $T_i$ 换成 $x_i$ 所得的元素 $f(x_i)$。事实上，对仿射空间 $k'^I$ 的每个固定点 $x$，立即就能核实 $\varepsilon_x$ 满足 $k$-代数同态的各条性质。

反过来，任意 $u\in\Hom_{k\text{-alg}}(P_I,k')$ 都确定 $k'^I$ 的一个点，也就是说，$u$ 都具有 $\varepsilon_x$ 的形式。

\Section{4}{坐标属于 $k$-代数的解集 $V_S(k')$}
\textbf{提要。}寻找多项式方程组 $S$ 的 $k$-值解，是一个“太困难”的算术问题，我们暂且将它搁置。这里为凯勒和格罗滕迪克所研究的较容易的代数几何问题，即在任意 $k$-代数中寻找解，提供一些理由。

考虑一族多项式 $(f_j)_{j\in J}$，其中 $J$ 是另一个任意固定的指标集。我们首先会想到，对“多项式方程组 $S$”
\[
S:\quad f_j(x)=0\qquad(j\in J)\tag{4.1}
\]
作通常的理解：研究这些方程的解，就是寻找点 $x\in k^I$，使每个 $j\in J$ 都满足 $\widetilde f_j(x)=0$。这种观点源自 $\mathbb R$、$\mathbb C$ 上仿射平面的经典情形。\footnote{例如，设 $S$ 由实欧氏平面中一条直线 $L$ 和一个圆 $C$ 的方程组成。如果 $L$ 与 $C$ 不相交，我们会“接受” $\mathbb C^2$ 中的两个复数解。经典代数几何从一开始，就被交点问题从 $\mathbb R$ 引向了 $\mathbb C$。}这样固定 $k$ 后，就会出现

% PDF 28
\sourcepage{24}{28}
一系列可能十分深刻、难以解决的问题。例如：
\begin{enumerate}[label=\arabic*)]
\item 方程组 $S$ 有没有解？
\item 如果有解，解的总数是有限还是无限？
\item 如果有限，能否对解的个数给出精确或近似的估计？
\item 如果无限，在加上某些不等式后，解的个数是否会变成有限？当参数趋于某些极限值时，能否给出这些估计的渐近公式？例如，在平面中，只找以原点为中心、各边平行于坐标轴的某个正方形内的解；然后问，当正方形边长趋于无穷大时，解的个数有什么渐近估计。
\end{enumerate}
对这些问题的回答，会随环 $k$ 的性质而相差极大。例如，比较 $\mathbb R$ 与 $\mathbb C$：$\mathbb C$ 代数闭而 $\mathbb R$ 不代数闭，这就带来很大的差别。此外，若 $k$ 是有限域，或者取 $k$ 为整数环，所得答案的类型也会大不相同。

这些问题是最困难的问题；这就是算术。我们将采用稍有不同的方法：允许解不仅在 $k^I$，即“$I$ 型标准仿射空间”中变动，而且在任意 $k'^I$ 中变动，其中 $k'$ 是任意 $k$-代数。于是，把这个代数方程组随 $k'$ 变化的解集 $V_S(k')$，看作关于变量 $k'$ 的函子 $k'\mapsto V_S(k')$。粗略地说，仿射代数几何就是研究这类函子及与之密切相关的

% PDF 29
\sourcepage{25}{29}
函子。

从实系数代数方程组的实数解过渡到复数解，这一经典技巧，可以看作向这个方向迈出的第一步：取 $k'=\mathbb R$、$\mathbb C=\mathbb R[1,i]$ 等，参见原书第23页脚注。然而，经典几何学家总是想着固定某个域：起初是 $\mathbb C$，后来是任意代数闭域；只在需要补救时才接受变动的“定义域”。凯勒，后来尤其是格罗滕迪克，提出在任意 $k$-代数中描述解；若没有函子语言，表述起来将极其繁琐。

\Section{5}{描述方程组解的函子 $V_S:\G_k\to\Sets$}
\textbf{提要。}虽然我们往往希望摆脱仿射嵌入，仍先定义从 $k$-代数到集合的函子 $\mathbb E^I$，即 $I$ 型仿射线性空间：它给每个 $k'\in\Ob\G_k$ 配上坐标属于 $k'$ 的仿射空间 $\mathbb E^I(k')=k'^I$，即 $k'\mapsto k'^I$。

设 $S$ 是任意一个 $(I,J)$ 型多项式方程组，系数属于基环 $k$。解函子 $V_S$（也就是“由 $S$ 定义的代数簇”）可以先作为 $\mathbb E^I$ 的子函子引入：$V_S\hookrightarrow\mathbb E^I$。它给每个 $k$-代数 $k'$ 配上 $S$ 在仿射空间 $\mathbb E^I(k')$ 中的解集。后面再摆脱 $\mathbb E^I$。\footnote{译注：原书提要中“由……定义的代数簇”的字母误排为 $k$，这里按语意写为 $S$。}

用格罗滕迪克本人的话说：

对于变动的 $k$-代数 $k'$（$k'\in\Ob\G_k$，参见第1节），我们考虑仿射空间 $k'^I$。希望把 $k'^I$ 解释为从 $P_I$ 到 $k'$ 的 $k$-代数同态的集合，也就是 $\Hom_{\G_k}(P_I,k')$，参见第2节：
\[
k'^I\xrightarrow{\sim}\Hom_{k\text{-alg}}(P_I,k')=\mathbb E^I(k').\tag{5.1}
\]
这个双射关于 $k'$ 是函子性的。于是，由 $k'\mapsto\Hom_{\G_k}(P_I,k')$ 定义的函子
\[
\mathbb E^I:\G_k\longrightarrow\Sets\tag{5.2}
\]
称为 $k$ 上的 $I$ 型标准仿射空间。
